\documentclass[12pt]{article}
\usepackage[utf8]{inputenc}
\usepackage[T1]{fontenc}
\usepackage[margin=1in]{geometry}
\usepackage{setspace}
\usepackage{xcolor}
\usepackage[colorlinks=true, urlcolor=blue!60!black, linkcolor=blue!60!black]{hyperref}
\setlength{\parindent}{2em}
\setlength{\parskip}{1.25em}
\renewcommand{\baselinestretch}{1.0}
\begin{document}
\begin{center}
{\LARGE \textit{Reed v. Town of Gilbert}, 576 U.S. 155 (2015)\par}
\vspace{0.8em}
{\large Case Brief --- Four-Page Edition\par}
\end{center}
\vspace{1em}

\section{Memory Jogger}

A town sign code gave a church's temporary directional signs harsher size and time limits than ideological or political signs; the Court held that a law that facially distinguishes signs by the topics or messages they convey is content based and triggers strict scrutiny regardless of benign motive, and struck the code (\href{https://supreme.justia.com/cases/federal/us/576/155/}{576 U.S. 155}).

\section{Detailed Case Facts}

Gilbert, Arizona's Sign Code banned outdoor signs without a permit but exempted 23 categories. Ideological signs could be 20 square feet with no time limit. Political signs could be 16 to 32 square feet for 60 days before a primary through 15 days after a general election. Temporary directional signs for a meeting of a religious, charitable, or nonprofit group were limited to 6 square feet, four per property, and 12 hours before to 1 hour after the event.

Good News Community Church, a small congregation without a building, and its pastor, Clyde Reed, posted 15 to 20 signs each Saturday announcing Sunday services and removed them around midday Sunday. The town twice cited the Church for exceeding the 12-hour window and omitting the event date, and confiscated a sign.

\section{Procedural History}

The District of Arizona denied a preliminary injunction, and the Ninth Circuit affirmed, holding that a ``cursory examination'' of a sign did not make the Code content based. On remand the district court granted summary judgment to Gilbert, and the Ninth Circuit again affirmed under \textit{Hill v.\ Colorado}, finding the categories content neutral for lack of censorial motive. The Supreme Court reversed on June 18, 2015.

\section{Judicial Votes}

Reversed, \textbf{9--0} in the judgment, with a \textbf{6--3} majority opinion.
\begin{itemize}
\item \textbf{Majority:} Justice Thomas, joined by Chief Justice Roberts and Justices Scalia, Kennedy, Alito, and Sotomayor.
\item \textbf{Concurrences:} Justice Alito, joined by Justices Kennedy and Sotomayor; Justice Breyer (in the judgment); Justice Kagan, joined by Justices Ginsburg and Breyer (in the judgment).
\item \textbf{Dissents:} None.
\end{itemize}

\section{Holding}

A law that distinguishes speech on its face by its communicative content, whether by subject matter, function, or purpose, is content based and triggers strict scrutiny regardless of benign motive, content-neutral justification, or lack of viewpoint animus. Gilbert's limits on temporary directional signs were hopelessly underinclusive as to aesthetics and traffic safety, so the Code failed strict scrutiny and violated the First Amendment.

\section{Analysis of Opinions}

\textbf{Thomas (majority).} Two tracks trigger strict scrutiny: laws content based on their face, and facially neutral laws that cannot be ``justified without reference to the content of the regulated speech'' under \textit{Ward v.\ Rock Against Racism}. Gilbert's Code was ``paradigmatic'' facial discrimination, since a sign's size and duration turned on whether it quoted John Locke, endorsed a candidate, or announced a meeting. Benign motive is irrelevant at step one because future officials can weaponize content-based laws, and \textit{Ward} applies only after a law is found facially neutral. Subject-matter discrimination is barred even without viewpoint discrimination, and speaker- or event-based distinctions reflecting content preferences are no different. Assuming the town's interests compelling, the Code was ``hopelessly underinclusive.'' Towns retain content-neutral tools: size, materials, lighting, moving parts, portability, and public-property bans.

\textbf{Alito (concurring).} Content-neutral rules on size, location, lighting, electronic displays, zones, on- versus off-premises signs, signs per mile, and time limits on signs for a ``one-time event'' remain valid, as does government speech.

\textbf{Breyer (concurring in the judgment).} Because regulation inevitably distinguishes by content, content discrimination should be a ``rule of thumb'' within a proportionality inquiry.

\textbf{Kagan (concurring in the judgment).} Automatic strict scrutiny threatens sensible ordinances, from ``Hidden Driveway'' warnings to historical markers. Gilbert's rules failed ``even the laugh test,'' so the Court need not have announced a sweeping rule.

\textbf{Later developments.} \textit{Barr v.\ AAPC} (2020) struck a robocall exception; \textit{City of Austin} (2022) held on- and off-premises distinctions content neutral; \textit{Chiles v.\ Salazar} (2026) relied on \textit{Reed} for a conversion-therapy ban.

\section{Examples: Future Applications}

\begin{enumerate}
\item \textbf{Same-Side (Invalid): Topic-Based Lawn Signs.} Oakridge allows only real estate, civic festival, and garage sale signs, each with its own limits. Classifying by subject is content based, and a real estate sign clutters as much as a bake-sale sign.
\item \textbf{Same-Side (Invalid): School-Zone Picketing.} A ban on picketing near schools exempts labor disputes. As in \textit{Police Department of Chicago v.\ Mosley}, the exemption turns on subject matter.
\item \textbf{Opposite-Side (Valid): Revised Neutral Ordinance.} Oakridge limits all temporary signs by size and days, bans off-premises digital billboards, and bars signs on utility poles. These are content neutral under \textit{Reed} and \textit{City of Austin}.
\item \textbf{Opposite-Side (Outside \textit{Reed}): Town Banners.} A town's own holiday banners are government speech, which Alito's concurrence preserved.
\item \textbf{Fence-Sitter: Uniform Event-Sign Limits.} A city limits signs for any one-time event to seven days before it. Alito listed such limits as permissible, but officials must read the sign, a tension Kagan identified.
\end{enumerate}

\section{Critique}

\textbf{Overbreadth of strict formalism.} As Kagan and Breyer emphasized, treating every facial subject-matter classification as presumptively unconstitutional either imperils benign administrative and safety rules or pressures courts to dilute strict scrutiny.

\textbf{Thomas versus Alito.} Kagan's footnote 1 noted that Alito's concurrence, needed for the majority, called one-time event limits content neutral, while Thomas struck Gilbert's rule for singling out signs about an event's time and place. \textit{City of Austin} largely resolved the tension in favor of Alito's line.

\section{Key Quotations}

\begin{quote}\sloppy
``Content-based laws---those that target speech based on its communicative content---are presumptively unconstitutional and may be justified only if the government proves that they are narrowly tailored to serve compelling state interests.''
\end{quote}
\noindent --- Justice Thomas, Opinion of the Court, 576 U.S.\ at 163.

\begin{quote}
``A law that is content based on its face is subject to strict scrutiny regardless of the government's benign motive, content-neutral justification, or lack of `animus toward the ideas contained' in the regulated speech.''
\end{quote}
\noindent --- Justice Thomas, Opinion of the Court, 576 U.S.\ at 165.

\begin{quote}
``Thus, a speech regulation targeted at specific subject matter is content based even if it does not discriminate among viewpoints within that subject matter.''
\end{quote}
\noindent --- Justice Thomas, Opinion of the Court, 576 U.S.\ at 169.

\end{document}
